\section[title={2024-09-10 - Vierfeldertafel},reference={2024-09-10-Vierfeldertafel}]

\startframedtask
S. 248

\stopframedtask

\textpartdivider

{\bf Beispiel am Urnenmodell}

In einer Urne sind {\bf 10} Kugeln, {\bf 5} davon sind Markiert
(Ereignis $M$). Also $P(M)=\frac{5}{10}=50\%$. Es gibt allerding {\bf
drei} von {\bf vier} roten Kugeln, welche Markiert sind und und {\bf
zwei} von {\bf sechs} nicht rote Kugeln. Wenn man nun beim ziehen vorher
schon weiß, welche Farbe die Kugel hat, bevor man die Markierung sieht,
verändert sich die Wahrscheinlichkeit auf $\frac{3}{4}=75\%$.

Man bezeichnet die Wahrscheinlichkeit für eine Markierung ($M$) unter
der Bedingung rot ($R$) als {\bf bedingte Wahrscheinlichkeit} und
schreibt.

\margintext{Das bedingende Ereignis $R$ wird als Index notiert. Man liest $P_R(M)$:
\quotation{Wahrscheinlichkeit von $M$ unter der Bedingung $R$}
}

\startformula 
P_R(M)=\frac{3}{4}=75%
 \stopformula

\startframedtask
{\bf Satz:} $P_E(F)$ ist die {\bf bedingte Wahrscheinlichkeit}

\stopframedtask

\section[title={Aufgaben},reference={aufgaben}]

Eine Münze wird dreimal geworfen. Berechnen Sie $P_E(F)$ und $P_F(E)$.
Beschreiben Sie die gesuchten Wahrscheinlichkeiten in Worten.

\startframedtask
\startitemize[a,packed][stopper=)]
\item
  E: \quotation{Beim zweiten Wurf lag \quote{Zahl} oben.} F:
  \quotation{Es lag dreimal \quote{Zahl} oben}
\stopitemize

\stopframedtask

$P_E(F)$ verändert sich zu \quotation{Es lag zwei mal Zahl oben}, da der
Zweite Wurf bereits Zahl hat. Also ergibt sich die Wahrscheinlichkeit
$P_E(F)=\frac{1}{4}$ anstatt $P(F)=\frac{1}{6}$

$P_F(E)$ da dreimal \quote{Zahl} oben liegt, dann liegt auch beim
zweiten Wurf \quote{Zahl} oben. Also eine Wahrscheinlichkeit von
$P_F(E)=1$

\startframedtask
\startitemize[a,packed][start=2,stopper=)]
\item
  E: \quotation{Beim ersten Wurf lag \quote{Zahl} oben} F: \quotation{Es
  lag genau einam \quote{Zahl} oben}
\stopitemize

\stopframedtask

$P_E(F) = \frac{1}{3}$: Möglichkeiten: $100,110,111$ (1: Zahl, 0: nicht
Zahl)

$P_F(E) = \frac{1}{3}$: Möglichkeiten: $100,010,001$

\startframedtask
{\bf Was wir wissen müssen}

\startitemize[packed]
\item
  Vierfeldertafeln aus und in sachzusammenhänge
\item
  Baumdiagramme hin und her
\item
  Was bedeutet bedingte Wahrscheinlichkeit
\item
  Fachsprache
\item
  Formeln kennen
\stopitemize

\stopframedtask
